\documentclass[11pt,a4paper]{article}
\usepackage[T1]{fontenc}
\usepackage[utf8]{inputenc}
\usepackage[margin=23mm]{geometry}
\usepackage{mathpazo,amsmath,amssymb,bm}
\usepackage{graphicx,booktabs,longtable,array,tabularx,multirow}
\usepackage{caption,subcaption,float,placeins}
\usepackage[numbers,sort&compress]{natbib}
\usepackage{xcolor,hyperref,url,microtype}
\usepackage{listings}
\hypersetup{colorlinks=true,linkcolor=blue!45!black,citecolor=green!35!black,urlcolor=blue!60!black,pdftitle={Geometry-dependent nuclear-data sensitivity of tritium production with runtime OpenMC perturbations}}
\graphicspath{{figures/}}
\setlength{\parskip}{0.35em}
\setlength{\parindent}{0pt}
\setlength{\emergencystretch}{2em}
\captionsetup{font=small,labelfont=bf}
\lstset{basicstyle=\ttfamily\small,breaklines=true,frame=single,backgroundcolor=\color{black!3},columns=fullflexible}
\newcommand{\LiSix}{\ensuremath{{}^{6}\mathrm{Li}}}
\newcommand{\LiSeven}{\ensuremath{{}^{7}\mathrm{Li}}}
\newcommand{\BeNine}{\ensuremath{{}^{9}\mathrm{Be}}}
\newcommand{\F}{\ensuremath{{}^{19}\mathrm{F}}}
\newcommand{\E}{\mathbb{E}}
\newcommand{\Var}{\operatorname{Var}}
\newcommand{\Cov}{\operatorname{Cov}}
\newcommand{\dd}{\mathrm{d}}
\newcommand{\figone}[4]{\begin{figure}[tbp]\centering\includegraphics[width=#2\linewidth]{#1}\caption{#3}\label{#4}\end{figure}}
\title{Geometry-dependent nuclear-data sensitivity of tritium production:\\runtime OpenMC perturbations, Morris screening, polynomial chaos and a DEMO-inspired Li$_2$O tokamak}
\author{Research manuscript draft\\\small Author names and affiliations to be completed before submission}
\date{6 November 2026\\\small Evidence snapshot of the modified OpenMC repository; not a peer-reviewed publication}
\begin{document}
\maketitle
\begin{abstract}
A runtime nuclear-data perturbation workflow is documented and assessed for fusion tritium-production studies. Native OpenMC reaction cross sections and supported independent secondary-energy and angular tables are modified through C++/C/Python interfaces, with native sampling, validation, derived-cross-section rebuilding and exact rollback. A critical scoring correction replaces cached MT205 production arrays with evaluation-verified reaction-specific production tallies. Five saved benchmark studies---a lithium sphere, FLiBe sphere, PbLi sphere and a layered Li/Be cube in two irradiation orientations---are examined using global Morris screening followed by independently validated sparse Legendre polynomial chaos expansions (PCEs) and coefficient-based Sobol indices. Between 163 and 706 parameters were screened; 24--30 inputs were retained. Held-out total-response $R^2$ values range from 0.97925 to 0.99818. The same layered material inventory exhibits a change in the Be-9 elastic total-order index from 0.0913 for a layer-parallel beam to 0.3891 for a layer-normal beam, demonstrating the relevance of transport geometry and irradiation conditions. Li-7 breakup production, particularly the evaluated MT56 branch, Li-6 MT105, and geometry-specific Be, F and Pb scattering/multiplication channels emerge as priorities for further nuclear-data investigation. These are conditional rankings under hypothetical independent perturbations, not covariance-based measurement priorities: three Morris studies remain unconverged, PDF modes have smaller individual bounds than cross sections, and several selected PCE bases exclude interactions by construction. Finally, a full-360-degree native-CSG DEMO-inspired tokamak with a Li$_2$O/Be blanket is described. Its higher-statistics segmented demonstration gives verified Li+Be production of $1.1774\pm0.0082$ tritons per source neutron. Tokamak sensitivity rankings have not yet been calculated. All numerical claims are traceable to saved artifacts, and the portable manuscript contains equations, figures, complete selected-input tables and reproduction instructions.
\end{abstract}
\textbf{Keywords:} fusion neutronics; tritium breeding; nuclear data; OpenMC; Morris; polynomial chaos; Sobol indices; geometry dependence; Li$_2$O blanket.
\clearpage
\tableofcontents
\clearpage
\section{Scientific motivation and transport formulation}
\label{sec:physics}
A breeding blanket must convert fusion neutrons into recoverable tritium while extracting heat and protecting structures and magnets. Nuclear-data priorities therefore depend on the response of interest and on the neutron histories established by the blanket. A high cross section or a high nuclide atom fraction alone does not determine response sensitivity. Earlier HCPB mock-up work combined detailed experimental geometry, reaction-rate measurements and sensitivity/uncertainty calculations \citep{hcpbsensitivity}. The present work uses deliberately simple sensitivity priors to identify channels warranting more detailed evaluation, differential measurements or integral validation. It does not replace evaluated covariance processing or experimental benchmarking.

OpenMC provides continuous-energy Monte Carlo transport and constructive solid geometry (CSG) \citep{openmc}. The existing repository has been extended to perturb loaded nuclear data directly, without reprocessing an evaluation or regenerating HDF5 libraries for each trial. This report audits and synthesises the available saved calculations; no new Morris campaign or tokamak sensitivity campaign was executed for manuscript preparation. Literature context is distinct from the original, repository-specific numerical results.

\subsection{Neutron transport and production response}
For steady fixed-source transport, let $\psi(\bm r,E,\bm\Omega)$ denote angular flux per unit source strength. A compact continuous-energy equation is
\begin{align}
\mathcal L(\bm x;\mathcal G)\psi &= q,\label{eq:transport}\\
\mathcal L\psi &= \bm\Omega\cdot\nabla\psi+\Sigma_t(\bm r,E)\psi
-\int_0^\infty\!\int_{4\pi}\mathcal K(\bm r;E',\bm\Omega'\!\to E,\bm\Omega)\psi(\bm r,E',\bm\Omega')\,\dd\Omega'\dd E'.
\end{align}
Here $\mathcal G$ includes material boundaries, gaps and boundary conditions; $\bm x$ is the declared nuclear-data parameter vector. The kernel includes reaction selection, secondary multiplicity and conditional energy-angle laws. Non-fission multiplication, such as $(n,2n)$, is included where present. Scalar flux is $\phi(\bm r,E)=\int_{4\pi}\psi\,\dd\Omega$.

For a specified set of verified triton-producing channels $\mathcal C$, the spatial-bin response is
\begin{equation}
R_{T,j}(\bm x)=\int_{V_j}\!\sum_n N_n(\bm r)\sum_{r\in\mathcal C_n}
\int_0^\infty y_{t,nr}(E)\sigma_{nr}(E;\bm x)\phi(\bm r,E;\bm x)\,\dd E\dd V.\label{eq:rate}
\end{equation}
$N_n$ is atom density, $\sigma$ microscopic cross section and $y_t$ tritons per reaction. OpenMC handles the barn-to-area conversion and target density in its nuclide-resolved tracklength scorer. In a source-normalised calculation, $R_{T,j}$ is tritons/source neutron. A score for a reaction is a triton yield only when its production multiplicity has been established. Sum $R_T=\sum_jR_{T,j}$; divide by $V_j$ for a bin-average production density per source neutron. Multiplication by an independently specified neutron source rate gives production per second. These units are not interchangeable.

For one DT neutron per fusion, an idealised neutronic breeding ratio is $\mathrm{TBR}=R_T$. An externally irradiated small sphere's production per source neutron should not be interpreted as a plant fuel-cycle breeding performance. Plant self-sufficiency also requires fuel processing, inventories, losses, availability and realistic penetrations.

\subsection{Why sensitivity depends on geometry}
At fixed source and geometry, differentiating Eq.~\eqref{eq:transport} gives
\begin{equation}
\mathcal L\frac{\partial\psi}{\partial x_i}=-\frac{\partial\mathcal L}{\partial x_i}\psi,
\qquad R_T=\langle r_T,\psi\rangle.
\end{equation}
With adjoint importance $\mathcal L^\dagger\psi^\dagger=r_T$, the formal first derivative is
\begin{equation}
\frac{\partial R_T}{\partial x_i}=\left\langle\frac{\partial r_T}{\partial x_i},\psi\right\rangle
-\left\langle\psi^\dagger,\frac{\partial\mathcal L}{\partial x_i}\psi\right\rangle.\label{eq:adjoint}
\end{equation}
The first term is a direct response-cross-section contribution; the second contains transport, spectrum, multiplication and leakage changes. Both forward flux and adjoint importance depend on geometry, source and materials. This equation is explanatory: the saved workflow uses finite perturbations and Monte Carlo, not an implemented adjoint solver.

For one perturbed production channel, the exact finite-change identity at fixed atom density is
\begin{equation}
\Delta R=\int N\left[\Delta\sigma\,\phi_0+\sigma_0\Delta\phi+\Delta\sigma\Delta\phi\right]\dd E\dd V.\label{eq:decomp}
\end{equation}
Thus an increase in elastic scattering can increase residence time and Li-6 capture, but also reduce flux above a Li-7 breakup threshold. The sign and net magnitude are not fixed by the elastic reaction's abundance. The existing constituent tallies support mechanism hypotheses, but the benchmark archives do not provide a complete energy-space decomposition of Eq.~\eqref{eq:decomp}; an exact direct/indirect numerical separation is not claimed.

\section{Runtime modifications to OpenMC and response verification}
\label{sec:runtime}
\subsection{Native data path and scope of the changes}
The relevant software stack is the existing modified branch \texttt{xs-change}, based on Git revision \texttt{fd99f66a20fd4841f66ce73eafac653d9ef8d2f2}, with uncommitted changes. Saved provenance identifies OpenMC \texttt{0.1.dev16025+gfd99f66a2.d20260924}. A revision alone does not uniquely identify this dirty working tree; per-study hashes and manifests are important. These changes belong to this research fork and are not asserted to be standard upstream OpenMC interfaces.

The data path is evaluated data $\to$ processed HDF5 $\to$ native C++ reaction/distribution objects $\to$ C interface $\to$ owned Python snapshots/candidates $\to$ validated native mutation $\to$ standard OpenMC sampling. The configured library is ENDF/B-VIII.0 \citep{endfb}. No evaluated file is rewritten in the sampling loop. Inspection uses the loaded objects, so a plotted supported angular or outgoing-energy distribution can be read back from the same representation sampled by transport.

\begin{table}[tbp]\centering\small
\caption{Key research-fork facilities, established from source inspection. Public Python names are exact; some lower-level legacy XS functions use a native nuclide index.}
\begin{tabularx}{\linewidth}{p{0.22\linewidth}X}
\toprule Facility & Implementation and purpose\\\midrule
XS inspection & \texttt{get\_nuclide\_xs}, \texttt{nuclide\_energy\_grid}, threshold/MT metadata getters in \texttt{openmc.lib}; native pointwise reaction arrays.\\
XS replacement/rebuild & \texttt{new\_nuclide\_xs\_with\_threshold}, \texttt{nuclide\_rebuild\_derived\_xs}; threshold-preserving replacement and derived-array consistency.\\
Secondary inspection & \texttt{get\_reaction\_products}, \texttt{get\_emission\_distribution}, \texttt{get\_angular\_distribution}; representation, frame, grid, PDF/CDF and interpolation metadata.\\
Native secondary changes & \texttt{set\_emission\_spectrum}, \texttt{set\_angular\_distribution}; supported existing independent tabular laws on preserved grids.\\
Native sampling checks & \texttt{sample\_emission\_energy}, \texttt{sample\_angular\_distribution}; call native samplers, not Python replacements.\\
Combined rollback & \texttt{nuclear\_data\_perturbation(xs=..., energy=..., angle=...)}; candidate validation, native deep copies, mutation, rebuild, rollback.\\
\bottomrule
\end{tabularx}\end{table}

The native implementation resides principally in \texttt{src/nuclide.cpp}, \texttt{src/secondary\_data.cpp}, the angle/energy distribution classes and C API declarations. Python bindings and orchestration reside in \texttt{openmc/lib/core.py}, \texttt{secondary.py} and \texttt{perturbation.py}. The shared \texttt{analysis\_codes/} package holds discovery, candidate generation, runner, screening, PCE and Sobol logic; model/config directories remain separate. Source anchors are listed in Appendix~\ref{app:source}.

The high-level transactional workflow is more restrictive than inspection. It admits independent uncorrelated continuous-tabular outgoing-energy tables and tabular angular laws. Isotropic, analytic, discrete-only and correlated representations can be inspectable without being safely mutable by this adapter. CorrelatedAngleEnergy, Kalbach--Mann and N-body phase space are not silently converted to independent PDFs. Fission, unresolved-resonance, multipole or other unsupported XS states are rejected by the native backup/validation machinery as applicable. Their absence from parameter rankings is not evidence of zero physical importance. Preserving PDF support and kinematic bounds also does not make independently perturbed shapes a new self-consistent evaluated reaction model: recoil/charged-particle energy balance and joint-law covariance require additional physics review.

\subsection{Validation, rebuilding and rollback}
Candidates are finite and nonnegative, preserve thresholds, grids, frame and reaction identity, retain discrete masses, and satisfy native PDF/CDF consistency. The context first acquires owned native backups, then applies changes while transport is not processing a batch. For perturbed nuclides, \texttt{Nuclide::rebuild\_derived\_xs()} clears and reconstructs the derived transport array from nonredundant reactions. The inspected implementation updates total, disappearance/absorption, fission, neutron-fission and photon-production contributions. A native consistency check follows. This is not a guarantee that every independently stored redundant MT is regenerated.

Snapshots are deep copies; writable pointers into native arrays are not exposed. Exit, including exceptions, restores the acquired native state in reverse order. This is rollback atomicity with serial access, not a database-style concurrent transaction. The runner verifies baseline and candidate readback and records failed evaluations instead of reusing contaminated state. Repeated evaluations use a live library session where supported, resetting simulation/tally state and seed between trials; worker processes keep separate OpenMC state. No independent transport kernel or alternative secondary sampler is introduced.

\subsection{Cached MT205 and the corrected production response}
\label{sec:mt205}
MT205 denotes $(n,Xt)$ production. In the processed library it is a redundant production cross section stored as reaction data. Physically one expects
\begin{equation}
\sigma_{205}(E)=\sum_r y_{t,r}(E)\sigma_r(E).\label{eq:mt205}
\end{equation}
However, the modified transport rebuild does not reconstruct the cached MT205 reaction array after changing its contributing partial reactions. The generic tally scorer interpolates the stored MT205 array if that score is requested. Consequently, with a modified partial cross section, a cached-MT205 score responds through the changed flux but misses the direct production-cross-section change. Schematically,
\begin{equation}
R_{205}^{\rm stale}(\bm x)=\int N\sigma_{205,0}\phi_{\bm x}\,\dd E\dd V,
\quad
R_T^{\rm correct}(\bm x)=\int N\sum_r y_{t,r}\sigma_{r,\bm x}\phi_{\bm x}\,\dd E\dd V.
\end{equation}
The baseline equality of the two responses does not validate their perturbed equality. An earlier cached-MT205 study therefore cannot be used as a complete production-sensitivity result; the corrected datasets analysed here use reaction-specific scores.

The response adapter verifies original evaluation metadata, processed-file hashes and baseline cross-section reconstruction. For the verified files its mapping is
\begin{align}
R_{T,6}&=R_{\LiSix,105},&
R_{T,7}&=\sum_{m=52}^{82}R_{\LiSeven,m},\\
R_{T,\mathrm{Be}}&=R_{\BeNine,700}+R_{\BeNine,701},&
R_{T,\mathrm{F}}&=R_{\F,105}.\label{eq:mapping}
\end{align}
Each included channel has verified unit triton yield. Li-7 MT51 is nonproducing; MT52--82 have the evaluation's LR=33 breakup flag. Li-7 MT33 is absent. Be-9 MT105 is a redundant parent of the scored MT700/701 channels and must not be counted twice. The PbLi response is explicitly lithium production: Pb nuclei remain in transport and perturbation discovery, but unverified Pb production channels are not assumed zero or invented. The tokamak response has its own Li+Be scope.

MT56 is an evaluated Li-7 inelastic branch with breakup metadata, not simply harmless elastic redistribution. Its saved metadata give $Q_M=-2.4667$ MeV and $Q_I=-4.6300$ MeV. The latter is a branch energy-loss parameter, not the laboratory threshold itself; native threshold-grid metadata are used for screening. Its importance is therefore physically compatible with a direct tritium contribution. Experimental proposals should address total breakup yields and differential spectra, not assume that a specific evaluated bin is a universally separable physical level.

\subsection{Correlated score uncertainty and controlled checks}
For batch $b$ with $H$ histories, a tracklength production estimator is
\begin{equation}
\widehat R_{T,b}=\frac{1}{H}\sum_{h\in b}\sum_{\ell\in h} w_\ell\ell\sum_{nr}N_n y_{t,nr}(E_\ell)\sigma_{nr}(E_\ell).
\end{equation}
Channel scores must be summed in each batch before error estimation:
\begin{equation}
\overline R_T=\frac{1}{B}\sum_b\widehat R_{T,b},\quad
s_{\overline R_T}^2=\frac{\sum_b(\widehat R_{T,b}-\overline R_T)^2}{B(B-1)}.
\label{eq:batch}
\end{equation}
Adding channel variances alone omits $2\sum_{r<s}\Cov(\overline R_r,\overline R_s)$. The implementation retains this covariance and uses common-seed paired batches for controlled perturbation diagnostics where appropriate.

Saved lithium verification gives baseline 1.172084, +20\% Li-6 MT105 response 1.199946, and +20\% Li-7 MT56 response 1.201743 tritons/source. Increasing Li-7 elastic by 20\% gives 1.157033: its Li-6 contribution rises by 0.018991 while its Li-7 contribution falls by 0.034042. The total decrease is thus consistent with competing transport pathways. Exact restoration returns the same saved baseline under controlled seeds. These are diagnostic means, not new uncertainty intervals; the production verification CSV does not itself provide confidence intervals for all paired differences. In FLiBe the analogous elastic perturbation increases the total slightly, further illustrating model dependence. These verification baselines use separate histories and should not be confused with the independent high-fidelity baseline values in the results table.

\section{Perturbation probability model and statistical methods}
\label{sec:methods}
\subsection{Independent sensitivity priors and native-grid constraints}
Each retained scalar coefficient is mapped from a unit coordinate $u_i\sim U(0,1)$ to
\begin{equation}
\delta_i=a(2u_i-1),\qquad a=0.20,\qquad \Var(\delta_i)=a^2/3.\label{eq:prior}
\end{equation}
These independent priors are hypothetical controlled sensitivity ranges, not evaluated covariance uncertainty. A symmetric prior need not produce a symmetric output distribution, equal numbers of increases/decreases in output, or zero mean response shift. Nonlinearity, normalisation and transport change the push-forward distribution.

For an XS whole-reaction parameter, the analysed operation is
\begin{equation}
\sigma'_r(E)=\sigma_{r,0}(E)(1+\delta_r).\label{eq:xs}
\end{equation}
Threshold zeros and native energy grids are preserved. The toolkit also has positive exponential parameterisations, but Eq.~\eqref{eq:xs} defines the reported whole-reaction studies. A native threshold above the 14 MeV source maximum removes \emph{that reaction's perturbation parameters}, including associated energy/angular modes; it does not delete the reaction from transport or from the production score. This rule is specific to the bounded source/non-fission studies, not a universal cutoff for arbitrary multiplying or broad-spectrum systems.

For tabulated PDFs, multiplying every ordinate by 1.2 and renormalising would make no shape change. Instead the shared strict-bound generator constructs nodal energy modes from $z=E_{\rm out}/\max(E_{\rm out,max},1\,\mathrm{eV})$, with $f_1=z$, $f_2=z^2$, and angular modes from $z=\mu=\cos\theta$,
\begin{equation}
P_1(\mu)=\mu,\qquad P_2(\mu)=(3\mu^2-1)/2.
\end{equation}
Let $\mathcal I_h$ be integration using the original table's interpolation, and $p_0$ its continuous portion. Strict mode centres and bounds each nodal basis:
\begin{equation}
\bar f_k=\frac{\mathcal I_h[p_0 f_k]}{\mathcal I_h[p_0]},\quad
 g_k=\frac{f_k-\bar f_k}{\max_j|f_k(x_j)-\bar f_k|},\quad
\widetilde p_j=p_{0,j}\left(1+\frac12\sum_{k=1}^{2}\delta_k g_{k,j}\right).\label{eq:shape}
\end{equation}
A zero basis is handled numerically without dividing by zero. A final normalisation enforces the unchanged continuous mass $M_c=1-\sum_{j=1}^{n_d}P_j$:
\begin{equation}
p'_j=\widetilde p_j\frac{M_c}{\mathcal I_h[\widetilde p]},\qquad
\sum_{j=1}^{n_d}P_j+\mathcal I_h[p']=1.\label{eq:pdfnorm}
\end{equation}
Discrete probabilities and support remain fixed. Provided the baseline and candidate satisfy tolerance, strict mode verifies the final pointwise ratio against $1\pm a$. The CDF is rebuilt from native segment areas, not independently guessed. Positivity follows from $|g_k|\le1$ and the budget bound, with additional finite-value checks.

\textbf{Unequal individual amplitudes matter.} The factor $1/2$ in Eq.~\eqref{eq:shape} remains even if one mode is not selected. One scalar shape coefficient varied alone can change the unnormalised local PDF by at most $\pm10\%$, whereas a scalar XS coefficient changes its XS by $\pm20\%$. Two shape modes jointly share the $\pm20\%$ budget. Consequently the reported family fractions are valid for this declared input model, but do not establish an intrinsic XS-versus-angle-versus-energy hierarchy independent of prior choice. A first-order variance contribution scales approximately with the square of the amplitude. Equalising coefficient bounds is not the same as equalising physical data perturbations or evaluated uncertainties.

\subsection{Interpolation, discrete mass and native sampling}
For adjacent native points $(x_j,p_j)$, segment areas are
\begin{equation}
A_j=\begin{cases}
(x_{j+1}-x_j)p_j,&\text{histogram},\\
\frac12(x_{j+1}-x_j)(p_j+p_{j+1}),&\text{linear--linear}.
\end{cases}\label{eq:areas}
\end{equation}
The supported mutating adapter does not silently approximate other interpolation rules. CDF increments are $A_j$; discrete-line CDF entries use cumulative probability masses. A histogram must be plotted as a step, not as a linearly interpolated PDF. Angular PDFs are defined in $\mu$ on $[-1,1]$, with reference frame preserved. If expressed as density in degrees,
\begin{equation}
p_{\theta_{\deg}}(\theta_{\deg})=p_\mu(\cos\theta)\sin\theta\,\pi/180.
\end{equation}
Relabelling an axis without that Jacobian would not give a density per degree. This report uses native $p(\mu)$.

An uncorrelated law factorises its native energy and angle components. A correlated law instead has structure such as
\begin{equation}
p(E_{\rm out},\mu\mid E_{\rm in})=p(E_{\rm out}\mid E_{\rm in})p(\mu\mid E_{\rm in},E_{\rm out}),
\end{equation}
which must not be replaced by an invented independent angular PDF. The independent-tabular studies exclude these unsupported mutations and retain their reasons in discovery inventories. Native sampler tests and saved readback plots support the implemented path, but unsupported channels require separate future development.

\subsection{Morris elementary-effects screening}
Morris screening \citep{morris} explores the full eligible parameter set efficiently before selecting the expensive surrogate dimension. With $d$ parameters, $r$ trajectories and $L=6$ grid levels,
\begin{equation}
\Delta=\frac{L}{2(L-1)}=0.6,\qquad N_{\rm Morris}=r(d+1).\label{eq:morriscost}
\end{equation}
Each random trajectory changes each coordinate once, with valid positive or negative grid steps. The elementary effect in the \emph{unit-coordinate} convention is
\begin{equation}
EE_{i,t}=\frac{R(\bm u_{t,k+1})-R(\bm u_{t,k})}{u_{i,t,k+1}-u_{i,t,k}}.\label{eq:ee}
\end{equation}
The physical coefficient step is $2a\Delta=0.24$, not 0.6 and not a one-sided 20\% step. The derivative with respect to $\delta$ differs by a factor $2a=0.4$. Across trajectories,
\begin{equation}
\mu_i=\frac1r\sum_tEE_{i,t},\quad
\mu_i^*=\frac1r\sum_t|EE_{i,t}|,\quad
s_i=\left[\frac1{r-1}\sum_t(EE_{i,t}-\mu_i)^2\right]^{1/2}.\label{eq:morrisstats}
\end{equation}
$\mu$ is signed, $\mu^*$ avoids cancellation, and $s_i$ reflects nonlinear/interaction variation \emph{plus Monte Carlo noise}. It is not a pure interaction measure.

For independent endpoint estimates the EE variance is $(s_A^2+s_B^2)/\Delta^2$. For same-seed endpoints,
\begin{equation}
\widehat{\Var}(EE)=\frac{s_A^2+s_B^2-2\widehat{\Cov}(\widehat R_A,\widehat R_B)}{\Delta^2}.\label{eq:eenoise}
\end{equation}
The workflow uses paired batch differences where configured, rather than assuming common seeds always reduce variance. Whole-trajectory bootstrap resampling retains dependence among adjacent effects. The saved convergence diagnostics compare split-half rank correlation, top-ten overlap, relative confidence widths and statistical resolution of leading effects. A resolved effect or a good PCE does not by itself establish Morris convergence; a small unresolved effect cannot be declared zero.

Selection is performed within XS, energy and angle families, with reserved diagnostic inputs. Thus the retained set is not simply the global top five, nor a proof that excluded parameters are negligible. The complete retained parameter identities and all screened Morris statistics are supplied with the manuscript.

\subsection{Independent training design and Legendre PCE}
After screening, the workflow runs a \emph{new} high-history Latin-hypercube design over the selected coordinates. Training and held-out validation use distinct seeds. Morris trajectory observations are \textbf{not automatically pooled} into these training data: they are low-history, path-structured and vary coordinates that may be excluded from the reduced surrogate. They could only be reused after defining a compatible model and noise treatment, which is not the reported procedure.

For $\xi_i=2u_i-1\sim U(-1,1)$, the probability-orthonormal polynomial is
\begin{equation}
\Psi_{\bm\alpha}(\bm\xi)=\prod_{i=1}^{d_a}\sqrt{2\alpha_i+1}\,P_{\alpha_i}(\xi_i),\qquad
\E[\Psi_{\bm\alpha}\Psi_{\bm\beta}]=\delta_{\bm\alpha\bm\beta}.
\end{equation}
The expansion is \citep{sudret,blatman}
\begin{equation}
\widehat R(\bm\xi)=\sum_{\bm\alpha\in\mathcal A}c_{\bm\alpha}\Psi_{\bm\alpha}(\bm\xi).\label{eq:pce}
\end{equation}
Total-degree bases have $\binom{d_a+p}{p}$ candidate terms. The implemented hyperbolic set obeys
\begin{equation}
\mathcal A=\left\{\bm\alpha:\left(\sum_i\alpha_i^q\right)^{1/q}\le p,\quad
|\operatorname{supp}\bm\alpha|\le2\right\},\qquad q=0.75.\label{eq:basis}
\end{equation}
Candidate orders are two and three. At $p=2$, a pair term $(1,1)$ has hyperbolic degree $2^{1/0.75}\approx2.520>2$: \textbf{all interaction terms are excluded}. At $p=3$, pure pair terms can enter, but three-way interactions remain prohibited. This is an important interpretation limit of the saved models.

The NumPy backend uses inverse-error weighted ridge fitting followed by weak-coefficient pruning/refitting, while retaining univariate main effects. Its threshold is relative to the largest fitted nonconstant coefficient; it is not the LARS algorithm. With design matrix $A_{k\alpha}=\Psi_\alpha(\bm\xi_k)$ and weights $W_{kk}=1/\max(s_k,s_{\rm floor})$, the ridge stage solves
\begin{equation}
\widehat{\bm c}=\arg\min_{\bm c}\{\|W(A\bm c-\bm y)\|_2^2+\lambda\bm c^\top D\bm c\},\quad D=\operatorname{diag}(0,1,\ldots,1),\quad\lambda=10^{-6}.\label{eq:wls}
\end{equation}
Writing the resulting linear estimator as $B\bm y$, its conditional observation-noise covariance is $B\operatorname{diag}(s_k^2)B^\top$. This excludes model discrepancy, basis-selection uncertainty and evaluated nuclear-data covariance. Four-fold training-only cross-validation chooses the polynomial order; the final independent validation set is not used for that selection.

The optional OpenTURNS backend \citep{openturns} provides weighted LARS/corrected-leave-one-out selection adapted to the same orthonormal basis and coefficient diagnostics. The installed adapter prewhitens observations and basis explicitly; its weight coordinate is not a physical uncertain input. The reported studies use the \textbf{NumPy backend}. No OpenTURNS accuracy advantage or nuclear-physics correction is inferred from its availability. A backend comparison must reuse identical data and held-out tests and cannot fix stale scoring, missing channels or an unsuitable prior.

\subsection{Validation, response distributions and Sobol decomposition}
For $n_v$ held-out observations,
\begin{equation}
R^2=1-\frac{\sum_k(\widehat R_k-R_k)^2}{\sum_k(R_k-\overline R)^2},\qquad
\mathrm{RMSE}=\sqrt{n_v^{-1}\sum_k(\widehat R_k-R_k)^2}.\label{eq:val}
\end{equation}
The saved ``normalised RMSE'' divides by $|\overline R|$, not the input-induced standard deviation. Both scales are considered here. Residuals must also be compared with transport standard errors. A high $R^2$ can coexist with residuals larger than Monte Carlo noise and can fail to resolve a weak individual Sobol component.

For an orthonormal PCE,
\begin{equation}
\E[\widehat R]=c_{\bm0},\qquad V=\Var(\widehat R)=\sum_{\bm\alpha\ne\bm0}c_{\bm\alpha}^2.\label{eq:pcevar}
\end{equation}
The variance attribution under independent inputs \citep{sobol,sudret} is
\begin{equation}
S_i=\frac{\sum_{\operatorname{supp}\bm\alpha=\{i\}}c_{\bm\alpha}^2}{V},\qquad
S_{T_i}=\frac{\sum_{\alpha_i>0}c_{\bm\alpha}^2}{V}.\label{eq:sobol}
\end{equation}
For a family/group $G$, its closed and total indices are
\begin{equation}
S_G=\frac{\sum_{\varnothing\ne\operatorname{supp}\bm\alpha\subseteq G}c_{\bm\alpha}^2}{V},\quad
S_{T_G}=\frac{\sum_{\operatorname{supp}\bm\alpha\cap G\ne\varnothing}c_{\bm\alpha}^2}{V}.\label{eq:groups}
\end{equation}
Summing individual total indices would double-count interactions and is not how the implementation groups coefficients. The saved resolution flags use component signal relative to propagated coefficient noise. They are not comprehensive confidence intervals on the indices. In particular, $S_T-S=0$ for an additive basis by construction and cannot exclude an unmodelled physical interaction.

The surrogate is sampled at 100,000 independent input vectors for response histograms and quantiles. Those cheap evaluations do not create additional transport information or reduce model bias. Direct validation observations are overlaid as finite samples with transport noise, not as an exact reference PDF. All distributions and indices are conditional on the retained inputs and assumed independent priors, with other inputs fixed at baseline.

\section{Benchmark models, archived settings and computational effort}
\label{sec:models}
Five current benchmark result directories contain the corrected production studies. These are the available completed analyses; older prototype Li$_2$O/UQ directories referred to in development history are not present in this checkout and are not represented as additional validated datasets. The separate Li$_2$O tokamak demonstration is discussed last in Section~\ref{sec:tokamak}.

The authoritative evidence is the saved configuration, source XML, model metadata, response mapping and checkpoint outputs, rather than an assumption that every current \texttt{config.py} exactly describes every historical run. Copies and SHA-256 hashes of the report's input artifacts are supplied. All benchmark sources are monoenergetic 14 MeV external pencil beams, with vacuum outer boundaries, natural lithium and the indicated material temperatures. They are diagnostic models, not complete power-plant blankets.

\begin{table}[tbp]\centering\small
\caption{Actual saved benchmark definitions. Geometrical sizes are target dimensions; a 1 cm surrounding vacuum region contains the external source. The Li-only density is a configured artificial value.}
\begin{tabularx}{\linewidth}{p{0.16\linewidth}p{0.17\linewidth}X p{0.17\linewidth}}
\toprule Model & Target & Materials and density [g/cm$^3$] & Beam\\\midrule
Li sphere & Radius 40 cm & Natural Li; 1.94; 900 K. Not realistic lithium-metal density. & From $z=-40.01$ cm, $+z$\\
FLiBe sphere & Radius 20 cm & Li$_2$BeF$_4$, Li:Be:F=2:1:4 atom ratio; 1.94; 900 K & From $z=-20.01$ cm, $+z$\\
PbLi sphere & Radius 40 cm & Li 15.7 atom\%, natural Pb 84.3 atom\%; 9.4; 900 K & From $z=-40.01$ cm, $+z$\\
Layers, parallel & Cube 20 cm per side & Alternating 1 cm Li (0.534) and Be (1.85) layers along $x$; 293.6 K & $(0.5,0,-10.01)$ cm, $+z$\\
Layers, perpendicular & Same cube and layers & Same atom inventories and material definitions as parallel case & $(-10.01,0,0)$ cm, $+x$\\
\bottomrule
\end{tabularx}\end{table}

\subsection{Li-only sphere}
This model isolates Li-6/Li-7 competition, production and scattering in a relatively optically thick, artificially dense lithium medium. It is useful as a numerical sensitivity benchmark but its 1.94 g/cm$^3$ density must be corrected before interpreting it as a physical lithium-metal experiment. Its high Li-7 production and substantial return into Li-6-producing energies need not transfer to a thinner or realistic-density breeder region. No comparison with a density-0.534 run at otherwise identical conditions has been made here.

\subsection{FLiBe and PbLi spheres}
FLiBe introduces Be multiplication/scattering and fluorine spectral/production pathways. The verified response includes Li, Be and F, not just lithium. PbLi introduces natural lead isotope scattering and $(n,2n)$ transport effects, while the specified response remains lithium-produced tritium. The different radii, compositions, densities and response scopes prevent assigning the difference between these sphere results solely to chemistry or solely to geometry. PbLi's atom fractions are the actual implemented 15.7/84.3 values; they should not be silently replaced by a nominal 17/83 composition.

\subsection{Layered Li/Be cube and a controlled irradiation change}
Both layered cases contain ten lithium slabs and ten beryllium slabs. The geometry and inventory are unchanged; the source direction and entrance position change. In the parallel case the initial beam travels through a central lithium slab. In the perpendicular case it encounters alternating media, entering lithium at the upstream face. Beam normalisation remains one source neutron. The comparison therefore probes the coupling of geometry and irradiation, not a pure material-content change and not an independently isolated shape variation at identical source conditions. Subsequent scatter and return paths can move neutrons between layers in both cases.

\figone{layers_geometry.png}{0.80}{Layered cube diagnostic from the saved model. Alternating Li/Be material regions are explicit native CSG. The perpendicular study retains these same layers and changes the beam direction.}{fig:layersgeom}

\subsection{Screening/training statistics and cost}
All saved studies use 20 Morris trajectories on six levels, then 512 independent training and 96 independent validation points plus a high-history baseline. The screening histories per evaluation are 100,000 for Li and 10,000 for the other cases. High-history evaluations use 1,000,000 histories for Li/PbLi and 100,000 for FLiBe and the two cubes, with ten batches each. ``Low/high fidelity'' here denotes histories/statistical fidelity of the same model, not a separate approximate transport solver.

\begin{table}[tbp]\centering\small
\caption{Saved executed cost and screening outcome. Counts include small baseline/diagnostic overhead above $20(d+1)$. Histories are cumulative source histories, not collision counts.}
\resizebox{\linewidth}{!}{\begin{tabular}{lrrrrl}
\toprule Study & Screened $d$ & Removed modes & Morris histories & High histories & Stopped converged?\\\midrule
Li & 163 & 57 & 328,400,000 & 609,000,000 & yes\\
FLiBe & 246 & 59 & 49,440,000 & 60,900,000 & no\\
PbLi & 706 & 70 & 141,440,000 & 609,000,000 & no\\
Layers parallel & 171 & 59 & 34,440,000 & 60,900,000 & yes\\
Layers perpendicular & 171 & 59 & 34,440,000 & 60,900,000 & no\\
\bottomrule
\end{tabular}}\end{table}

Threshold filtering removes 19, 21, 32, 21 and 21 reaction targets, respectively, with all associated eligible modes. It is distinct from exclusions due to unsupported representations. These counts refer to parameters and native thresholds, not to a decision that all remaining data are equally influential or that all nuclear-data types have adapters. Every removed-reaction list is included in the portable data snapshot.

The recorded Morris wall times are approximately 707, 531, 3609, 261 and 317 s, and high-phase wall times 5533, 811, 12,976, 793 and 284 s in the above order. Worker counts differ, notably eight high-phase workers for the perpendicular cube versus one for the parallel cube; wall times are not a clean transport-throughput comparison. Resumable deterministic checkpoints retain seeds, parameter vectors, errors and restoration status. No new high-history transport was needed to prepare this report.

\section{Sensitivity and uncertainty results}
\label{sec:results}
\subsection{Morris screening: robust signals and remaining uncertainty}
Figure~\ref{fig:morris} compares leading elementary effects, while the complete searchable CSV tables contain all 163--706 screened parameters per model. Effects retain their unit-coordinate scaling, so raw magnitudes should not be compared across very different response levels without considering the scale of $R_T$.

\figone{morris_comparison.pdf}{1.0}{Ten largest total-response Morris effects per benchmark, with trajectory-bootstrap 95\% intervals. The restricted figure is a comparative view, not the full parameter inventory; complete statistics and resolution flags are supplied in \texttt{data/}.}{fig:morris}
\figone{morris_convergence.pdf}{1.0}{Saved ranking-stability diagnostics as trajectories accumulate. Dashed lines mark rank/overlap thresholds, but the full stopping criterion also checks effect resolution and interval widths. A satisfactory curve in one panel does not imply complete convergence.}{fig:morrisconv}

Li-7 MT56 and Li-6 MT105 lead the Li/FLiBe screening, PbLi is led by Li-6 MT105, and both layered models show a strong Li-7 MT56 signal. At 20 trajectories the final split-half rank correlations are 0.9778 (Li), 0.9208 (FLiBe), 0.9600 (PbLi), 0.9282 (parallel layers) and 0.8926 (perpendicular layers). The saved stopping criterion passes only Li and parallel layers. The three failures must remain explicit even though their leading signals and total-response PCE fits are useful. The present evidence does not establish a converged complete ordering of all weak parameters.

\subsection{PCE validation and probability distributions}
\begin{table}[tbp]\centering\small
\caption{Selected sparse-PCE fits for the total response. $d_a$ is the retained dimension; $p$ the selected order; $M$ the candidate hyperbolic basis size. The terms may have zero fitted coefficients after sparse pruning.}
\input{tables/validation.tex}\label{tab:validation}
\end{table}
The selected orders are three for Li/PbLi and two for FLiBe/both cubes. Held-out RMSE relative to the validation response standard deviation is approximately 4.24\%, 7.30\%, 4.36\%, 5.79\% and 14.33\%, respectively. These are more informative for sensitivity than error normalised only to a large nonzero mean. All total responses pass the saved acceptance rule, which permits a sufficiently high $R^2$ or noise-consistent residuals, with resolved surrogate variance.

The average squared noise-standardised residual is about 1.43 (Li), 1.21 (FLiBe), 9.69 (PbLi), 3.09 (parallel layers) and 9.40 (perpendicular layers). Thus high $R^2$ does not mean residual discrepancies are explained by transport noise, especially for PbLi and perpendicular layers. Model-form/truncation error remains and may matter to small sensitivity components. Component-resolution flags propagate observation noise only; they do not cover all sparse-model selection uncertainty.

\figone{layer_validation.pdf}{1.0}{Independent validation of the two layered studies, with one transport standard error on the native OpenMC horizontal axis and the ideal line. The perpendicular case has appreciably larger relative surrogate error.}{fig:layerval}

\begin{table}[tbp]\centering\scriptsize
\caption{Baseline and conditional PCE response distributions, in verified tritons/source neutron. Baseline errors are one batch standard error; PCE spread is input-induced variation, not a standard error. The 5--95\% range is not an evaluated nuclear-data confidence interval.}
\resizebox{\linewidth}{!}{\input{tables/distributions.tex}}\label{tab:distribution}
\end{table}
\figone{response_comparison.pdf}{1.0}{Baseline response and conditional input-induced coefficient of variation. Different geometry/material/response definitions prevent interpreting this as a ranking of breeder technologies.}{fig:responsecomparison}

The conditional coefficients of variation are about 2.38\% for Li, 4.74\% for FLiBe, 7.90\% for PbLi, 4.26\% for parallel layers and 4.84\% for perpendicular layers. They arise from the declared independent ranges and down-selected modes. They cannot be substituted for evaluated covariance uncertainties or directly compared with experimental error bars.

\subsection{Lithium benchmark: production competition and MT56}
\begin{table}[tbp]\centering\caption{Leading conditional Sobol indices: Li sphere. Full selected-input table in Appendix~\ref{app:indices}.}\input{tables/li_analysis_top.tex}\end{table}
Li-6 MT105 has $S_T=0.4251$ and Li-7 MT56 $S_T=0.4084$. The fact that Morris places MT56 narrowly first while the conditional PCE puts MT105 first is not contradictory: Morris measures finite-path response changes in the full screened space, whereas PCE variance attribution is squared and conditional on selected inputs. Near-tied rankings should not be overinterpreted.

Li-7 elastic XS contributes 0.0695 and its second angular mode 0.04395. Elastic scattering is consequential but not automatically dominant. MT56 directly contributes tritons, is open at 14 MeV (native threshold 5.29564 MeV), and can compete with scattering/slowing pathways. The controlled test in Section~\ref{sec:mt205} shows elastic-induced Li-6 enhancement and Li-7 reduction. Li-7 neutron-induced breakup can produce a triton while leaving a secondary neutron, unlike the neutron-absorbing Li-6 MT105 channel; event counts and triton yields therefore need not be bounded by one per source. The model's artificial Li density and thick geometry favour particular histories; this is not a universal ranking for natural lithium.

\begin{figure}[tbp]\centering
\begin{subfigure}{.49\linewidth}\includegraphics[width=\linewidth]{li7_mt56_xs.png}\caption{Li-7 MT56 cross section.}\end{subfigure}\hfill
\begin{subfigure}{.49\linewidth}\includegraphics[width=\linewidth]{li6_mt105_xs.png}\caption{Li-6 MT105 cross section.}\end{subfigure}
\caption{Saved native-readback baseline and perturbed cross sections for leading lithium channels. These are whole-reaction perturbation examples, not energy-resolved sensitivity profiles or measured covariance bands.}\label{fig:xs}
\end{figure}
\begin{figure}[tbp]\centering
\begin{subfigure}{.49\linewidth}\includegraphics[width=\linewidth]{li7_mt2_angle.png}\caption{Li-7 elastic angular mode.}\end{subfigure}\hfill
\begin{subfigure}{.49\linewidth}\includegraphics[width=\linewidth]{li7_mt16_energy.png}\caption{Li-7 $(n,2n)$ outgoing-energy mode.}\end{subfigure}
\caption{Supported native angular/energy perturbation examples. The metadata record readback, reference frame, selected incident energy, interpolation and normalisation. The outgoing-energy mode is statistically unresolved in this study; its inclusion illustrates the available representation, not a claim of strong influence.}\label{fig:pdfs}
\end{figure}
For the displayed energy law, native baseline continuous-PDF normalisation is 0.999999934 and both perturbed readbacks integrate to unity within numerical precision. Such checks establish valid mutation, not experimental validation of the distribution.

\subsection{FLiBe: fluorine transport joins lithium production}
\begin{table}[tbp]\centering\caption{Leading conditional Sobol indices: FLiBe sphere.}\input{tables/flibe_analysis_top.tex}\end{table}
Li-7 MT56 leads at 0.39294, followed by F-19 elastic at 0.20376, Li-6 MT105 at 0.18869 and F-19 MT22 at 0.12720. The Be-9 elastic index is 0.05096. Fluorine is therefore a transport-relevant constituent even when attention initially focuses on lithium. MT22 denotes neutron-plus-alpha emission and changes neutron energy and competition; it is not itself synonymous with a unit-triton channel. F-19 MT105 is a separately verified production contribution.

The baseline response separates into Li-6 0.046413, Li-7 0.141612, Be-9 0.003206 and F-19 0.008807 tritons/source. Its conditional PCE mean is 0.200048, with a 5--95\% interval 0.184463--0.215696. The Morris screening remains unconverged. The selected order-two hyperbolic basis is additive: equality of $S_i$ and $S_{T_i}$ in this result is a basis constraint, not a measurement of absent cross-family interactions.

\figone{flibe_analysis_tritium_distribution.png}{0.85}{Saved conditional FLiBe production distribution from PCE sampling. The independent transport observations are finite noisy validation data, not an exact reference density.}{fig:flibedist}

\subsection{PbLi: production, lead scattering and multiplication}
\begin{table}[tbp]\centering\caption{Leading conditional Sobol indices: PbLi sphere, lithium-production response only.}\input{tables/lipb_analysis_top.tex}\end{table}
Li-6 MT105 contributes $S_T=0.64080$. Pb-208 elastic and $(n,2n)$ contribute 0.11459 and 0.10659, followed by Li-7 MT56 at 0.05195. Lead isotope content affects scattering, multiplication and the emergent spectrum; the response itself counts only the verified lithium channels. These results motivate lead nuclear-data investigation for PbLi blankets, without implying a lead-production yield map has been established.

The baseline is 0.053483 tritons/source, and the conditional 5--95\% interval is 0.046774--0.060604. Although the total $R^2$ is 0.99808, residuals substantially exceed quoted transport noise, and Morris remains unconverged. The order-three basis permits pair terms; the Li-6 XS first/total separation is only about 0.00130 in this reduced model. Strong interactions outside that subspace, higher-order interactions and biases from excluded data are not ruled out.

\subsection{Layered cube: geometry/irradiation changes sensitivity ordering}
\label{sec:geometryresult}
\begin{figure}[tbp]\centering\includegraphics[width=\linewidth]{geometry_and_families.pdf}
\caption{Left: closed grouped family fractions under the declared priors and selected bases. Right: selected shared channels in the same layered cube under two beam orientations. The uncertainty model is reduced-space and the active sets are not identical; these are not universal nuclide importance factors.}\label{fig:geometrysobol}
\end{figure}
The parallel case is dominated by Li-7 MT56 ($S_T=0.79025$), Be-9 elastic (0.09125), Li-6 MT105 (0.02682) and Li-7 elastic (0.02204). The perpendicular case changes these to MT56 0.44806, Be elastic 0.38913 and Li-6 MT105 0.11178. Be elastic thus rises by a factor of about 4.26 in its \emph{conditional fraction}. The raw total-response Morris Be-elastic $\mu^*$ also rises from 0.009288 to 0.015886, while MT56 falls from 0.028783 to 0.017893. This supports a change beyond merely renormalising Sobol fractions.

At baseline the parallel source produces 0.216739 tritons/source; the perpendicular source gives 0.161106, a decrease of approximately 25.7\%. Li-7 production falls from 0.197670 to 0.118752, whereas Li-6 rises from 0.015604 to 0.027914 and Be-9 from 0.003464 to 0.014440. These scores are consistent with a redistribution of high-energy and slowed-down paths when the beam crosses interfaces. Their explanation is a transport hypothesis supported by constituent totals, not a completed energy-resolved reaction-history attribution.

The two cases retain 30 and 28 uncertain inputs, respectively, because family selection and support differ. A rigorously matched quantitative test should impose the union of the same uncertain inputs, identical priors and additional convergence checks in both orientations. Nevertheless, their shared inventory with changed irradiation and strongly different shared-channel effects demonstrates why nuclide content alone is insufficient. More generally, Eq.~\eqref{eq:adjoint} predicts changes with blanket thickness, penetrations, layer placement and region of interest even at fixed composition.

\subsection{Family contributions and limitations of weak-mode conclusions}
Cross-section closed-family fractions are 0.9496 (Li), 0.9905 (FLiBe), 0.9945 (PbLi), 0.9969 (parallel layers) and 0.9882 (perpendicular layers). Angular fractions are approximately 0.04995, 0.00932, 0.00542, 0.00302 and 0.01137. Energy fractions are all below $5\times10^{-4}$ in these selected models, and are marked unresolved except in the perpendicular case. This is evidence about \emph{supported independent tabular modes, their smaller per-mode amplitudes, the selected inputs and total-response observables}. It cannot justify disregarding correlated emission physics, other angular modes, spatial/energy-resolved responses or evaluated covariance.

\FloatBarrier
\section{Priorities for nuclear-data investigation and measurement}
\label{sec:priorities}
\subsection{What can legitimately be prioritised}
A large sensitivity index identifies a pathway under a specific response, geometry and input model. It does not show that the underlying evaluated datum is inaccurate or that a new measurement is more valuable than improving another datum. For a small perturbation about baseline,
\begin{equation}
\Var(R)\approx\bm g^\top C_\delta\bm g,\qquad g_i=\partial R/\partial\delta_i.
\label{eq:covariancepriority}
\end{equation}
The present independent uniform prior has diagonal variance $a^2/3$ before mode budgeting. A realistic measurement priority also requires evaluated covariance, existing experimental coverage, achievable precision, parameter correlations and experimental cost. If only independent variance $C_{ii}$ is reduced by a factor $\eta$, the linearised variance reduction is $(1-\eta)g_i^2C_{ii}$; correlation changes make that expression incomplete. This illustrates why equal-range sensitivity and value of information are different quantities. No covariance-based optimisation is claimed here.

The ordered list below is therefore a \textbf{provisional investigation sequence}, chosen from recurrence, maximum observed impact and model-specific physics. It is not a universal ordering or a claim of demonstrated deficiencies. Before commissioning new measurements, compare evaluated files and covariances with experimental coverage in EXFOR \citep{exfor} and fusion-oriented evaluation/benchmark work such as FENDL \citep{fendl}. A different intended blanket may reorder the list.

\begin{table}[tbp]\centering\small
\caption{Evidence-led investigation sequence under the reported hypothetical priors. Numerical fractions are conditional $S_T$, not evaluated uncertainty shares.}
\begin{tabularx}{\linewidth}{c p{0.23\linewidth}X}
\toprule Order & Nuclear-data target & Evidence and recommended focus\\\midrule
1 & Li-7 breakup production, especially evaluated MT56 and the MT52--82 set & MT56 $S_T=0.408$ in Li, 0.393 in FLiBe and 0.448--0.790 in layers; only 0.052 in PbLi. Verify branch sums, unit yields, evaluation partitioning and total/differential breakup data before interpreting a bin-specific priority.\\
2 & Li-6 $(n,t)$, MT105 & $S_T=0.425$ in Li and 0.641 in PbLi, with direct production impact in every breeder. Audit covariance/standards consistency and energy-dependent yield/cross-section evidence across the actual transported spectrum.\\
3 & Be-9 elastic, MT2; associated production/multiplication channels & Elastic $S_T$ changes from 0.091 to 0.389 with layer orientation; inspect differential scattering, $(n,2n)$ and triton channels without redundant-parent double counting.\\
4 & F-19 elastic MT2 and MT22 in FLiBe; production MT105 & FLiBe elastic and MT22 indices 0.204 and 0.127. Prioritise spectral redistribution and reaction competition in a salt-specific integral experiment, not only lithium measurements.\\
5 & Pb-208 elastic MT2 and $(n,2n)$ MT16, then other Pb isotopes & PbLi indices 0.115 and 0.107. Investigate isotope-resolved scattering/multiplication and emergent spectra for PbLi applications. Not relevant in a Pb-free blanket.\\
6 & Differential angular data: Li-7 elastic, Be-9 elastic and F-19 elastic & Li elastic angular P2 $S_T\simeq0.044$; Be/F angular effects emerge in their geometries. Study frames, angle-energy correlation and response-dependent angular observables.\\
7 & Supported outgoing-energy shapes and unsupported correlated laws & Small selected independent-mode contributions do not establish unimportance of full emission data. Resolve noise, compare equal physical budgets and validate conditional energy-angle kernels before assigning measurement priority.\\
\bottomrule
\end{tabularx}\end{table}

\subsection{Suggested differential and integral experimental programme}
For Li-7, an energy scan should cover the open breakup region sampled by the application, with particular attention to the MT56 native threshold near 5.30 MeV and the 14 MeV source region. These are physically motivated \emph{candidate} energies, not a measured energy-sensitivity profile. Inclusive triton yield, outgoing-neutron spectra and angular distributions can test whether an apparent branch-specific effect is robust under different evaluation decompositions. Measuring an inclusive physical channel is often more transferable than treating a discretised evaluated state as independently uncertain.

For Li-6, measure or reassess production cross sections against the relevant established standards and covariances before concluding that additional differential data are needed. Slowed-down neutrons can dominate Li-6 production even with a 14 MeV source, so a source-energy-only measurement would not automatically resolve the application uncertainty. The current whole-reaction study cannot identify the optimal incident-energy measurement interval.

For Be/F/Pb, elastic angular distributions, inelastic energy transfer and multiplication affect the forward and adjoint fields. An integral programme could compare matched Li/Be slabs with beam parallel and perpendicular to layers, keeping material inventories fixed, and measure constituent production together with leakage and neutron spectra. Spatially resolved tritium-production measurements and leakage spectra would help distinguish compensation between direct production and transport. The present results motivate such a design; they do not substitute for measured integral validation.

\subsection{Geometry-specific rather than universal nuclear-data priorities}
The recurrence of Li-7 MT56 across several simple fast-beam models is a credible reason to investigate that evaluation and response mapping, particularly given the corrected direct-production score. It is not sufficient to announce MT56 as the dominant uncertainty of a DEMO blanket. The tokamak has enriched Li-6, oxygen, steel, tungsten, distributed isotropic emission, longer return paths and magnet/vessel structures; its importance weighting will differ. Be-9 dominance in some HCPB integral literature \citep{hcpbsensitivity} and Li-7 dominance in these particular benchmarks need not conflict: geometry, material fractions, source, response and uncertainty model all matter.

A defensible next campaign would repeat full-space screening to convergence, test the excluded parameter set, match parameter priors across geometry comparisons, permit needed interaction terms, and then replace the hypothetical priors with evaluated covariance information. These are recommendations, not additional methods or results silently added to the present study.

\FloatBarrier
\section{Final application model: DEMO-inspired \texorpdfstring{Li$_2$O}{Li2O} tokamak}
\label{sec:tokamak}
\subsection{Engineering reference and deliberate departures}
The new \texttt{tokomak\_analysis/} directory contains a full-360-degree native OpenMC CSG model, separate from the preceding small benchmarks. Published EU-DEMO design-space and magnet studies motivate the machine scale \citep{demodesign,magnets}; HCPB and shielding work motivate breeder/multiplier/cooling concepts \citep{hcpb,shielding}; double-wall vessel structure is informed by containment studies \citep{vessel}. These sources describe evolving designs, not one geometrically complete common reference. The implemented dimensions combine documented values with explicitly assumed details.

\begin{table}[tbp]\centering\small
\caption{Tokamak baseline. Assumed detailed builds are not presented as published DEMO dimensions.}
\begin{tabularx}{\linewidth}{p{0.29\linewidth}p{0.24\linewidth}X}
\toprule Parameter & Implemented value & Status\\\midrule
Plasma major/minor radius & 9.10/2.90 m & Reference-inspired machine scale\\
Elongation / triangularity & 1.65 / 0.33 & Design-space shaping; faceted approximation\\
Fusion power / source & 2 GW / 14.1 MeV & Reference-scale power; isotropic plasma-volume source\\
TF / PF / CS modules & 16 / 6 / 6 & TF count reference-inspired; detailed placements/builds assumed\\
IB / OB blanket build & 75.5 / 98.2 cm & Reference radial-build motivation\\
Plasma--first-wall gap & 22.5 cm & Reference radial-build motivation\\
Tungsten armour & 0.2 cm & Representative concept value\\
FW steel/coolant/rear & 0.8 / 1.0 / 0.2 cm & Assumed split of 2 cm envelope\\
Breeder / multiplier & Li$_2$O / Be & Li$_2$O deliberately replaces reference HCPB ceramics\\
Li-6 enrichment / packing & 60 atom\% / 0.64 & Configurable; enrichment reference-inspired, packing assumed\\
Blanket modules & $16\times3\times16=768$ nominal & Assumed toroidal and IB/OB vertical-stack segmentation\\
Vessel inner/filler/outer & 3.6 / 50.4 / 6.0 cm & Simplified double-wall/filler split\\
\bottomrule
\end{tabularx}\end{table}

The reference HCPB ceramic blend is not Li$_2$O. The substitution changes lithium inventory, oxygen content, density, thermal properties and irradiation/purge behaviour. Li$_2$O solid density is 2.013 g/cm$^3$ and default pebble packing 0.64, with interstitial purge helium. Different breeder cycles can have different Li-6 enrichment. Be beds, cooling plates, EUROFER-like structures, manifolds, SS316L vessel, W armour, Cu/CuCrZr and homogenised Nb$_3$Sn winding packs are distinct regions/materials. This neutronics calculation does not establish Li$_2$O engineering qualification.

\subsection{Native CSG geometry and hierarchy}
The toroidal axis is $Z$; OpenMC lengths are cm. A faceted Miller-like plasma contour is generated from
\begin{equation}
R(\alpha)=R_0+a\cos[\alpha+\delta\sin\alpha],\qquad
Z(\alpha)=\kappa a\sin\alpha,\quad 0\le\alpha<2\pi,
\end{equation}
then revolved using installed native polygon/RZ helpers and analytic planes, cylinders and cones. An arbitrary elliptical-torus primitive is not assumed. Offset contours form the chamber/blanket/vessel builds. Repeated rotated universes supply the full toroidal assembly, with distinct D-shaped TF winding, insulation and case regions, individually placed PF coils, central supports/shielding and cased CS modules. The chamber has a lower opening connected to the divertor vacuum; there is no magnetic-equilibrium or true X-point calculation.

Blanket modules have toroidal gaps and sidewalls and separate inboard/outboard vertical stacks. Three breeding/multiplier/cooling cycles lie behind the first wall and ahead of backplate/manifold structures. Level 1 homogenises active interiors while retaining major layers/gaps; level 2 resolves radial beds; level 3 adds grouped helium slots, ribs and breeder-unit subdivisions. Explicit pebbles, pipe bends and complete manufacturing detail are not included. A lower divertor has two inclined W/CuCrZr/water/steel targets and a cooled supporting cassette, distinct from the blanket.

\figone{tokamak_xz.png}{0.90}{Native-cell-query poloidal slice of the DEMO-inspired tokamak. The source plasma, blanket, double-wall vessel, D-shaped TF envelopes, central column, PF coils and lower divertor are distinct CSG regions. Colours identify component/material regions; this is not an independently drawn transport geometry.}{fig:tokamakxz}
\begin{figure}[tbp]\centering
\begin{subfigure}{.49\linewidth}\includegraphics[width=\linewidth]{tokamak_blanket.png}\caption{Outboard blanket detail.}\end{subfigure}\hfill
\begin{subfigure}{.49\linewidth}\includegraphics[width=\linewidth]{tokamak_divertor.png}\caption{Lower divertor detail.}\end{subfigure}
\caption{Native-query engineering detail: repeated breeder/multiplier/cooling layers with structures and manifolds, and separately layered divertor targets. Thin features are limited by slice pixel size.}\label{fig:tokamakdetails}
\end{figure}

\subsection{Source normalisation, scores and material approximations}
The baseline source is isotropic, uniform in physical plasma volume. Native cylindrical proposals are weighted proportional to $R$ and reject outside the plasma cell; this correctly includes the cylindrical volume Jacobian. Native \texttt{TokamakSource} supports an optional radial profile, while an optional asymmetric source bank uses plasma-volume rejection solely for source generation. Nuclear secondary sampling remains native.

Assuming 17.6 MeV per DT fusion,
\begin{equation}
\dot N_n=\frac{P_f}{17.6\,\mathrm{MeV}}=7.092624\times10^{20}\ \mathrm{s}^{-1}\quad(P_f=2\,\mathrm{GW}),
\end{equation}
with unit conversion understood. At 14.1 MeV/neutron this is 1602.27 MW source-neutron power. A per-source heating score $h$ in eV/source corresponds to $P_h=\dot N_n h(1.602176634\times10^{-19})$ W. It is not obtained by multiplying by 2 GW directly.

The tokamak verified-production response includes Li-6 MT105, Li-7 MT52--82 and Be-9 MT700/701, with the same evaluation-pinning and baseline-reconstruction safeguards. Oxygen/structural triton-producing side channels have not been exhaustively mapped; this is explicitly a scoped Li+Be TBR. Per-module, poloidal-stack and inboard/outboard contributions are summed in each batch and checked against the global response. Flux spectra, mesh flux, boundary leakage, shield-interface current, $(n,2n)$ event rate, component heating and magnet/vessel exposure are saved. A single $(n,2n)$ score is not a full neutron balance.

Neutron-only calculations use \texttt{heating-local}; an additional small coupled neutron--photon demonstration uses collision-estimator \texttt{heating}. Temperatures and material compositions are approximations: neutron transport in magnet materials uses available evaluated temperatures near room temperature rather than a true 4.5 K data library; hot-water thermal scattering is not fully represented; representative winding/steel compositions are not certified engineering specifications. A Zr96 probability-table warning at 600 K in the supplied library is retained, not suppressed by changing data.

\subsection{Executed transport results and validation evidence}
\begin{table}[tbp]\centering\small
\caption{Executed native geometry-debug transport demonstrations. Production errors are one batch standard error. Different blanket fidelity changes transport, so agreement is not assumed.}
\begin{tabular}{lrrrr}
\toprule Fidelity/case & Histories & Unique cells & Li-only TBR & Verified Li+Be TBR\\\midrule
Global, level 1 & 10,000 & 220 & $1.22594\pm0.01758$ & $1.23261\pm0.01755$\\
Segmented, level 2 & 20,000 & 348 & $1.16626\pm0.00826$ & $1.17736\pm0.00823$\\
Detailed, level 3 & 5,000 & 524 & $1.13566\pm0.01822$ & $1.14700\pm0.01818$\\
Coupled smoke, level 1 & 400 & 220 & $1.20388\pm0.04269$ & $1.21087\pm0.04240$\\
\bottomrule
\end{tabular}\end{table}
The level-2 inboard and outboard contributions are $0.41333\pm0.00644$ and $0.76403\pm0.01012$ tritons/source. Their standard errors cannot be added in quadrature without their covariance. Its Li-6/Li-7/Be-9 means are 1.13007, 0.03619 and 0.01110. Enrichment and the reactor-scale spectrum produce a very different constituent balance from the natural-lithium fast-beam benchmarks, reinforcing the need for a separate tokamak sensitivity calculation rather than importing their rankings.

Construction takes roughly 1.0--1.1 s. Level-2 geometry has 13,216 physical leaf-cell instances, 2,049 surfaces and 17 used materials; detailed geometry has 21,664 leaf instances and 2,529 surfaces. The respective measured debug-run rates are approximately 753, 686 and 441 source histories/s for levels 1--3, including setup/verification/output overhead, with peak process memory roughly 0.85--0.95 GiB. These are not pure collision-kernel or large-machine scaling benchmarks.

A 20-test suite passed, covering composition, geometry construction, signed boundaries, XML, native surfaces, periodicity, source confinement, overlap/coverage sampling, module counts, native debug transport/conservation, VTK readback and CLI validation. All three final fidelities completed debug transport with no reported overlap/lost-particle failure. Nevertheless, 36 cells in the 20k run had fewer than ten overlap checks, so this is finite test evidence, not an exhaustive proof of nonoverlap. Exact interfaces and universe-local surface sharing resolved boundary problems found during development; difficult components were not deleted to obtain a clean run.

A shared 100,000-point native-query inventory comparison estimates Li differences relative to level 2 of $+0.188\%\pm1.954\%$ for level 1 and $-0.115\%\pm0.487\%$ for level 3 (paired standard errors). Estimates are consistent with similar inventories but do not prove exact inventory equivalence or sub-percent tolerance. Global homogenisation and lower truncation can have a real effect. Level-1 TBR differs visibly from level 2, and these demonstrations do not establish blanket-fidelity convergence.

\figone{tokamak_breeding.png}{0.94}{Level-2 poloidal-stack Li+Be production, with batch standard errors. Individual module totals are saved and checked against the global score.}{fig:tokamakbreeding}
\figone{tokamak_spectrum.png}{0.88}{Saved blanket neutron spectrum. Values are native energy-bin integrated volume-integrated flux, not a probability density or automatically per-eV flux. A precise energy-resolved sensitivity kernel has not been inferred from this plot.}{fig:tokamakspectrum}

\subsection{Three-dimensional export and outstanding research questions}
VTK export samples the authoritative CSG through native cell queries and maps physical instance paths to component, material, module, sector and inboard/outboard IDs. It provides a VTM multiblock assembly with VTU children, VTI images and VTP surfaces, without replacing the transport geometry. The full-reactor grid is 96$\times$96$\times$80, with voxels about 39.6$\times$39.6$\times$31.3 cm. Thin armour may disappear. A separate first-wall zoom uses 0.1$\times$0.75$\times$0.5 cm voxels. Both pass VTK readback. ParaView GUI inspection was unavailable; the supplied cutaway is rendered by Matplotlib from the same VTK surfaces because headless OpenGL failed.

\figone{tokamak_cutaway.png}{0.98}{Sampled three-dimensional CSG assembly cutaway. Coarse-voxel gaps/stepping in thin components are visualisation artefacts; this is not a boundary-conforming engineering mesh.}{fig:tokamak3d}

No tokamak Morris/PCE/Sobol campaign is available. No sensitive tokamak nuclide ordering, measured covariance uncertainty, converged deep-shielding result or self-sufficiency claim can therefore be reported. Zero magnet tallies at these histories are unresolved, not proof of sufficient shielding. Penetrations, realistic coolant routing, magnetic/structural self-consistency, neutron-data temperature treatment and Li$_2$O thermomechanics remain substantial limitations. This model supplies the next geometry on which the validated response/mutation machinery can be investigated; its TBR above unity is not a plant design conclusion.

\FloatBarrier
\section{Conclusions and scientific status}
The runtime framework makes controlled changes to the loaded native neutron data sampled by OpenMC. Its most important response correction is avoiding stale MT205 production arrays when partial cross sections change. Evaluation-verified partial reaction scores, per-batch summation and exact rollback establish a physically consistent basis for the reported sensitivity studies.

The saved benchmarks identify Li-7 breakup production and Li-6 MT105 as recurring leading pathways, with Be, F and Pb transport data important in the corresponding materials. The same layered inventory under two beam orientations changes Be-9 elastic conditional $S_T$ from about 0.091 to 0.389, and changes absolute elementary-effect magnitudes and constituent production rates. Geometry/irradiation and response choice must therefore accompany any statement of nuclear-data importance.

Sparse PCEs predict held-out total responses well at the declared ranges, but do not validate the full uncertainty model. Three screening runs remain unconverged, the selected PCE spaces differ, shape coefficients have smaller individual physical amplitudes than XS coefficients, and order-two hyperbolic bases exclude interactions by construction. These limits preclude a universal uncertainty ranking or a claim that weak energy effects and physical interactions are absent.

The ordered investigation sequence is Li-7 breakup/MT56, Li-6 production, geometry-relevant Be scattering/multiplication, FLiBe-specific F scattering/MT22, PbLi-specific Pb scattering/multiplication, then differential angular and unresolved/unsupported emission data. This is an evidence-led review list, not a prescribed experimental budget. Existing differential measurements, evaluated covariances and integral benchmarks must be reviewed before declaring data gaps or choosing new measurements.

The full-360-degree Li$_2$O tokamak is an explicit next application with completed CSG/debug-transport/VTK demonstrations. Its segmented verified Li+Be yield is $1.1774\pm0.0082$ tritons/source at 2 GW source normalisation, but its nuclear-data sensitivity ordering has not been calculated. Enriched breeder, structure, oxygen and distributed source make extrapolation of the simple-model rankings unjustified. The next defensible study is a geometry-qualified, converged sensitivity analysis of that model, followed by covariance-informed uncertainty assessment and targeted experimental comparison.

\begin{table}[tbp]\centering\small
\caption{Classification of conclusions and outstanding work.}
\begin{tabularx}{\linewidth}{p{0.26\linewidth}X}
\toprule Status & Findings\\\midrule
Numerically supported within declared studies & Corrected production responses; native mutation/readback/restoration; high held-out total-response $R^2$; strong selected channel effects; orientation-dependent shared-channel sensitivity.\\
Conditional or suggestive & Full-space importance ordering in unconverged studies; transport mechanism attribution from constituent totals; relative measurement priorities; weak family contributions.\\
Not established & Evaluated covariance uncertainty; complete all-constituent production for PbLi/tokamak; exact energy-space direct/indirect decomposition; absent physical interactions; geometry-independent ranking.\\
Demonstrated tokamak capabilities & Native CSG/XML/loading, three fidelity debug runs, scoped production/heating/flux scores, finite geometry tests, sampled VTK readback.\\
Not yet evaluated for tokamak & Morris/PCE/Sobol indices, energy-resolved sensitivity, high-statistics magnet shielding, engineering qualification and fuel-cycle self-sufficiency.\\
\bottomrule
\end{tabularx}\end{table}

\FloatBarrier
\bibliographystyle{unsrtnat}
\bibliography{references}
\clearpage
\appendix
\section{Source-code traceability}
\label{app:source}
Line numbers refer to the inspected working-tree snapshot; the distributed data manifest identifies copied artifacts. They can change after later edits. The report preparation script reads saved outputs and does not run transport.

\begin{longtable}{p{0.45\linewidth}p{0.48\linewidth}}
\caption{Audited implementation anchors.}\\
\toprule Source anchor & Scientific relevance\\\midrule\endfirsthead
\toprule Source anchor & Scientific relevance\\\midrule\endhead
\path{src/nuclide.cpp}:555, \texttt{Nuclide::rebuild\_derived\_xs} & Clears/rebuilds transport-derived arrays; does not reconstruct the cached MT205 reaction array.\\
\path{src/nuclide.cpp}:1496 & C wrapper for derived-XS rebuild.\\
\path{src/tallies/tally_scoring.cpp}:1071--1099 & Generic nuclide-resolved reaction scorer calls \texttt{get\_nuclide\_xs}, with density and flux weighting.\\
\path{openmc/data/neutron.py}:402--409,634--641 & Export/redundant handling preserves production MTs including 205 as stored redundant reactions.\\
\path{src/secondary_data.cpp}:598,615,668,697 & Reaction info, owned backups, restoration and XS validation C interfaces.\\
\path{src/secondary_data.cpp}:714,738,758,792 & Angular setter, native energy/angular sampling and snapshot creation.\\
\path{openmc/lib/secondary.py}:93,137,178,279 & Product/emission/angular inspection and native angle mutation wrappers.\\
\path{openmc/lib/perturbation.py}:68,203--313 & Native histogram/linear areas; prevalidated combined context, deep backups, apply/rebuild, reverse restoration.\\
\path{sandbox/analysis_codes/production.py} & Verified mapping, original evaluation evidence, baseline reconstruction, per-batch weighted response and controlled checks.\\
\path{sandbox/analysis_codes/discovery.py}, \path{thresholds.py} & Representation eligibility and threshold filtering; unsupported targets retained in audit.\\
\path{sandbox/analysis_codes/perturbations.py} & Strict centred, shared-budget nodal shape modes and native readback.\\
\path{sandbox/analysis_codes/morris.py} & Six-level trajectories, unit-coordinate effects, paired/noise statistics and bootstrap diagnostics.\\
\path{sandbox/analysis_codes/family_workflow.py}:62--129 & Independent high-history training/validation LHS designs and execution.\\
\path{sandbox/analysis_codes/pce.py}, \path{sobol.py} & Probability-orthonormal Legendre basis; weighted ridge/pruning; exact support-based variance decomposition.\\
\path{sandbox/analysis_codes/openturns_pce.py}, \path{validation.py} & Optional prewhitened LARS adapter; independent validation/noise diagnostics.\\
\path{tokomak_analysis/} & Separate parametric CSG application, verified channel tallies, native-query export and tests.\\
\bottomrule
\end{longtable}

\section{Complete selected-input Sobol results}
\label{app:indices}
Every retained parameter is listed below. The shorthand P1/P2 is an angular perturbation mode, E1/E2 an outgoing-energy mode; product/distribution identities remain in the supplied original CSV files. ``Resolved'' is the saved coefficient-noise diagnostic, conditional on basis and inputs; it is not a full uncertainty interval. A printed zero is a fitted coefficient result, not proof of physical irrelevance. Nonretained parameters have not been assigned zero indices.
\input{tables/li_analysis_sobol.tex}
\input{tables/flibe_analysis_sobol.tex}
\input{tables/lipb_analysis_sobol.tex}
\input{tables/be_li_layers_sobol.tex}
\input{tables/be_li_layers_x_beam_sobol.tex}

\section{Saved-data reproducibility and manuscript use}
The project contains \texttt{main.tex}, section sources, \texttt{references.bib}, figures, generated tables and a compact evidence snapshot. Upload the contents of the project ZIP to Overleaf, choose \textbf{pdfLaTeX} and \textbf{main.tex}, and compile. BibTeX is used, not Biber; standard TeX Live packages suffice. Author names, affiliations and acknowledgements should be completed by the research team. This is a technical manuscript draft, not an assertion of journal acceptance or experimental validation.

Build locally with:
\begin{lstlisting}
cd fusion_engineering_report
pdflatex -interaction=nonstopmode -halt-on-error main.tex
bibtex main
pdflatex -interaction=nonstopmode -halt-on-error main.tex
pdflatex -interaction=nonstopmode -halt-on-error main.tex
\end{lstlisting}

To regenerate comparative figures/tables from the repository's saved results without transport:
\begin{lstlisting}
cd /sandbox/openmc
/opt/openmc-venv/bin/python fusion_engineering_report/prepare_report.py
\end{lstlisting}
The packaged snapshot supports checking the quoted tables and archived observations outside the repository; regenerating from the current repository requires the original result directories. No nuclear-data libraries or binary OpenMC executable are included in the manuscript ZIP.

The framework's transport reproduction commands, subject to explicit configured budgets and installed matching library, are:
\begin{lstlisting}
python -m sandbox.analysis_codes Li_analysis --dry-run
python -m sandbox.analysis_codes Li_analysis --resume
python -m sandbox.analysis_codes FLiBe_analysis --resume
python -m sandbox.analysis_codes LiPb_analysis --resume
python -m sandbox.analysis_codes Be_Li_layers --resume
python -m sandbox.analysis_codes Be_Li_layers_x_beam --resume
# Refit saved data only, without new transport:
python -m sandbox.analysis_codes FLiBe_analysis --reanalyse
# Native-CSG tokamak demonstration:
python -m sandbox.tokomak_analysis.main --fidelity 2 \
 --particles 2000 --batches 10 --run --debug \
 --output sandbox/tokomak_analysis/results/reproduction
\end{lstlisting}
Saved checkpoints must be checked for matching configuration fingerprints before resumption. Editing model/prior/response definitions requires a new result set; old cached-MT205 values cannot be relabelled as corrected production. Current scripts or libraries may evolve, so preserve the manifest, code snapshot and actual per-run provenance when reproducing historical results.

\end{document}
